%% Research proposal only. No implementation or performance claim.
%% Compile separately from root.tex; keep out of a results submission until validated.
\documentclass[letterpaper,10pt,conference]{ieeeconf}
\IEEEoverridecommandlockouts
\overrideIEEEmargins
\usepackage{amsmath,amssymb,booktabs,url}
\title{Response-Aware Leaf Presentation through Bounded Probing and Regrasping}
\author{Research-method specification --- not an evaluated system}
\begin{document}
\maketitle
\begin{abstract}
This document specifies an unimplemented extension to an existing robotic leaf-inspection
system. Its candidate contribution is task-directed estimation and use of the local
motion response of an attached leaf, including an explicit option to abandon a poorly
coupled grasp without increasing the force budget. No performance results are claimed.
The plant model and both working robot pipelines remain unchanged.
\end{abstract}

\section{Scope and distinction}
Moving leaves to improve visibility is established~\cite{yao2025}, as are leaf grasping
for sensing~\cite{mortazavi2026} and online deformation mappings~\cite{hu2017}. The
candidate contribution is narrower: estimate whether a particular grasp transfers motion
to a selected sensing region, and use that estimate to choose bounded presentation,
regrasp, or release. An adaptive Jacobian by itself is not the novelty. This claim must
survive comparison with a standard adaptive visual servo and a simple low-response
regrasp rule. We estimate effective response, not Young's modulus or tissue health.

\section{Observation and identification}
Track two free-blade regions outside the contact envelope: one near the jaw and one
containing the desired sensing region. Estimate their positions and normals from RGB-D,
transform them to a common world frame, and retain correspondence across frames. Fit
quality, point count, occlusion and association consistency define a validity flag. Do
not predict the region by assuming it follows the gripper: that would assume the very
coupling being tested. A local tracker may use robot motion as a broad search proposal,
but must also admit a stationary leaf and reject uncertain associations.

Choose a fixed tangent basis $B\in\mathbb R^{3\times2}$ at each reference normal. For
small excursions, the regional normal feature is $z=B^\top(\mathbf n-\mathbf n_0)$.
Sign-align normals with the previous valid estimate; when identity or face orientation
is uncertain, abstain. Stack the two regional features in $y\in\mathbb R^4$ and use
the measured action increment $u=[\Delta h/\ell_0,\Delta q_g/\theta_0]^\top$, where
$h$ is lift and $q_g$ is actual tool roll. Fixed scales $\ell_0,\theta_0$ avoid mixed
units. Normal motion about its own axis is unobservable from normals alone; this
representation is not a full six-degree-of-freedom leaf pose.

After bilateral contact, execute a short sequence of small lift and reverse-roll probes
within the existing force and workspace limits. Probe amplitudes and dwell times are
selected on development trials and frozen before testing. Initial engineering candidates
are 2\,mm lift increments and $3^\circ$ roll increments, not biologically safe bounds.
Use separate lift and roll increments so the input matrix is sufficiently excited; wait
for a predefined settling criterion and subtract measured no-action drift. For valid
increments, fit the local response
\begin{equation}
 \Delta y_i = J u_i + b\Delta t_i + \epsilon_i,
\end{equation}
where $b$ captures residual drift. One regularized estimate is
\begin{equation}
 (\widehat J,\widehat b)=\arg\min_{J,b}
 \sum_i w_i\|\Delta y_i-Ju_i-b\Delta t_i\|^2
 +\lambda\|J\|_F^2.
\end{equation}
Compute prediction intervals from held-out development residuals, accounting for temporal
correlation; do not interpret a least-squares residual as a calibrated confidence bound.
Require adequate excitation and observation coverage before using the estimate. Force
is separately regulated by the existing inner loop and is not inferred from $J$.

\section{Response-aware action selection}
For a small candidate action set $\mathcal U$, predict the desired region's feature
change and its uncertainty. Let $e=y-y^\star$ and $W$ prioritize the sensing region.
Define the conservative predicted benefit
\begin{equation}
 L(u)=\|e\|_W^2-\|e+\widehat Ju\|_W^2
       -\beta\sigma_{\mathrm{pred}}(u)-\lambda_u\|u\|^2.
\end{equation}
Here $\sigma_{\mathrm{pred}}$ bounds uncertainty in the predicted cost reduction, in
the same cost units, and $\beta$ is fixed on development data. Reject actions outside
joint, corridor, roll, lift and force-monitor limits. Positive $L$ is a local prediction,
not a formal safety certificate; retain the raw-force stop and abort on unobserved contact.

If a valid action has positive conservative benefit, execute only one small increment,
observe the actual response, and update the model. Repeated prediction failures trigger
re-identification or release. If observations are invalid, first pause and reacquire the
region; do not equate missing data with a rigid or slipping leaf. If observations are
valid but no useful action remains, open, withdraw through the existing planner, and
regrasp at a different visible patch on the same leaf, subject to a fixed attempt budget.
Candidate patches must pass the existing reach and canopy checks. Their response is
unknown until probed; a measured response at one patch must not be assumed valid elsewhere.
Compare this decision with a matched random or distal regrasp baseline so that extra
attempts alone cannot explain an improvement.

A low regional response may reflect anchoring, bending, local buckling, finger
compliance, sliding or tracking error. Classify it as \emph{low transfer}, not slip.
Sliding requires an additional relative tangential displacement observation at the
contact or another independently validated sensor; force and normal change do not
uniquely identify the mode~\cite{zhou2022}.

\section{Termination and evaluation contract}
Terminate on measured leaf orientation error below a fixed tolerance, sufficient visible
region area, bilateral contact within the force window, and a continuous dwell. Example
development targets are a $15^\circ$ regional reorientation with $5^\circ$ tolerance and
2\,s dwell; register these before test runs and report achievable-goal coverage. Report
unreachable and unobserved goals in the denominator, not only completed grasps.

Keep the plant geometry, materials, FEM settings, attachments, damping and initialization
fixed. Compare on the same leaf, initial grasp patch, sensing region, goal, force cap and
time/attempt budget. The two hands form separate strata unless common contact geometry
and a common reachable subset support a paired comparison. Historical modulus sweeps are
context only, not part of the new experiment.

\begin{table}[h]
\centering\footnotesize
\caption{Required baselines and ablations; all are proposed, not measured.}
\begin{tabular}{@{}ll@{}}
\toprule
condition & question isolated \\
\midrule
Fixed lift/roll & current system \\
Camera-only active view & is contact needed for the sensing goal? \\
Fixed-gain visual servo & does leaf feedback suffice? \\
Adaptive visual servo, no regrasp & does grasp change matter? \\
Low-response heuristic regrasp & is the response model needed? \\
Response-aware, fixed grasp & separates adaptation from regrasp \\
Response-aware, no uncertainty gate & does abstention help? \\
Full method & combined task-directed policy \\
\bottomrule
\end{tabular}
\end{table}

Primary outcomes are paired presentation success and time to the first valid pose, with
timeouts included. Secondary outcomes include angular tracking error, sensing-region
availability, force exceedance duration, regrasp count and force exposure
$\int F_{\mathrm{sum}}\,dt$ (N\,s, not mechanical work). After release, measure residual
regional deformation and recovery time; neither proves absence of biological damage.
Report camera-estimator error against evaluation-only FEM trajectories without using
ground truth for control. Log both peak and terminal deformation and all invalid samples.
Use run/leaf-pair clustered analysis, not camera frames as independent trials. With only
one plant asset, uncertainty intervals concern that tested geometry, not a plant population.

\section{What would falsify the contribution?}
If fixed-gain feedback matches the full policy, the model is unnecessary. If heuristic
regrasp matches it, the candidate selection/estimation contribution is unsupported. If
camera-only viewing achieves equivalent sensing coverage without contact, manipulation
is not justified for that task subset. If the policy benefits only from a larger force or
attempt budget, the proposed claim fails. None of these comparisons has yet been run.

\bibliographystyle{IEEEtran}
\bibliography{refs}
\end{document}
